\section{Conclusion and Future Work}
\label{sec:conclusion}
%
% PURPOSE OF THIS CHAPTER
% Answer the research question directly, summarize demonstrated contributions,
% and identify extensions that remain outside the validated benchmark.
%
% \subsection{Answer to the research question}
% - State how a retained discrete-time generative mechanism supports controlled
%   datasets and multiple explanation ground truths.
% - Summarize what the experiments establish about recovery and reliability;
%   avoid claims not supported by Chapter 8.
%
% \subsection{Contributions}
% - Configurable expression-first time-series generation.
% - Multiple trajectories per mechanism and reproducible parallel batches.
% - Horizon-aware structural, signed, differential, and coalition-based reference
%   explanations actually implemented by the final version.
% - Model-independent adapter and metric interface.
% - Empirical comparison and operator-level failure analysis.
%
% \subsection{Limitations}
% - Fixed grammar and restricted innovation families.
% - Empirical filtering versus formal stability guarantees.
% - Computational cost and estimator uncertainty of coalition explanations.
% - Limits of interpreting synthetic-process ground truth as evidence about real
%   systems.
%
% \subsection{Future work}
% - More innovation distributions, heteroscedasticity, regime switching,
%   distributed memory, multiscale and richer multivariate mechanisms.
% - Stability certificates and more complete symbolic canonicalization.
% - Learned difficulty models and adaptive curriculum generation.
% - Pretraining priors that jointly predict the future, mechanism structure or
%   equation, and innovation distribution.
% - Explanation-aware training heads using generated attribution references, in
%   the general spirit of supervised explanation learning reviewed for
%   SHAPformer.
% - Additional symbolic/system-identification and explainable-by-design models.
